% Generated draft: requires mathematical and source review.
\documentclass[11pt]{article}
\usepackage[margin=25mm]{geometry}
\usepackage{fontspec}
\setmainfont{Latin Modern Roman}
\usepackage{amsmath,amssymb,mathtools,unicode-math}
\setmathfont{Latin Modern Math}
\usepackage{enumitem,needspace,xurl,hyperref}
\setlength{\parindent}{0pt}
\setlength{\parskip}{5pt}
\setlength{\emergencystretch}{3em}
\newcommand{\N}{\mathbb{N}}
\newcommand{\Z}{\mathbb{Z}}
\newcommand{\Q}{\mathbb{Q}}
\newcommand{\R}{\mathbb{R}}
\newcommand{\C}{\mathbb{C}}
\newcommand{\F}{\mathbb{F}}
\newcommand{\A}{\mathbb{A}}
\newcommand{\PP}{\mathbb{P}}
\newcommand{\E}{\mathbb{E}}
\newcommand{\eps}{\varepsilon}
\newcommand{\dd}{\,\mathrm{d}}
\begin{document}
\section*{UnsolvedMath ID: 30002158}
\textbf{Source problem code:} catalogue entry
\subsection{Definitions and mathematical statement}
The supplied record asks for a distribution of isomorphism classes of short exact sequences. It does not specify the domain of the field \(k\), the allowed primes \(p\), or the family, ordering, and counting convention used to define density. Consequently, the following gives a conditional formulation in the standard number-field setting; that setting is an explicit interpretive assumption, not a hypothesis supplied by the record.

Fix a number field \(k\) and a rational prime \(p\). Let \(\Z_p\) and \(\Q_p\) denote the ring of \(p\)-adic integers and its fraction field. For an elliptic curve \(E/k\), let \(E(k)\) be its group of \(k\)-rational points, and put
\[
E[p^\infty]=\bigcup_{n\geq 1}E[p^n](\overline{k}).
\]
Writing \(H^1\) for Galois cohomology and letting \(v\) range over all places of \(k\), define the Tate--Shafarevich group by
\[
\operatorname{Sha}(E/k)=\ker\!\left(H^1(k,E)\longrightarrow\prod_v H^1(k_v,E)\right).
\]
For an abelian group \(G\), write
\[
G[p^\infty]=\{g\in G: p^n g=0\text{ for some }n\geq 1\}.
\]
The \(p\)-primary Selmer group is
\[
\operatorname{Sel}_{p^\infty}(E/k)=\ker\!\left(H^1(k,E[p^\infty])\longrightarrow\prod_v H^1(k_v,E)[p^\infty]\right),
\]
where the maps are induced by localization and the Kummer sequence. Its natural short exact sequence of \(\Z_p\)-modules is
\[
\mathcal S_p(E/k):\qquad
0\longrightarrow E(k)\otimes_{\Z}(\Q_p/\Z_p)
\longrightarrow\operatorname{Sel}_{p^\infty}(E/k)
\longrightarrow\operatorname{Sha}(E/k)[p^\infty]
\longrightarrow 0.
\]
No finiteness assumption on \(\operatorname{Sha}(E/k)\) is added here.

Let
\[
S:\qquad 0\longrightarrow A\xrightarrow{i}B\xrightarrow{q}C\longrightarrow 0
\]
be a fixed abstract short exact sequence of \(\Z_p\)-modules. The assertion \(\mathcal S_p(E/k)\cong S\) means that there are \(\Z_p\)-linear isomorphisms between the corresponding three terms that commute with both arrows; it concerns the entire exact sequence, not merely the middle module.

\textbf{Distribution question.} Given \(S\), determine the density of elliptic curves \(E/k\) satisfying \(\mathcal S_p(E/k)\cong S\). More explicitly, if a counting convention consists of increasing finite nonempty sets \(\mathcal F_T\) of \(k\)-isomorphism classes of elliptic curves, exhausting a specified family as \(T\to\infty\), the corresponding unweighted question is to establish existence and determine the value of
\[
\delta_{k,p,\mathcal F}(S)
=\lim_{T\to\infty}
\frac{\#\{[E]\in\mathcal F_T:\mathcal S_p(E/k)\cong S\}}
{\#\mathcal F_T}.
\]
This formula makes one possible meaning of density explicit; the supplied record does not select these sets, require unweighted counting of isomorphism classes, or specify an alternative convention. A definitive formulation therefore requires the missing field and counting hypotheses. The supplied open-status label and distribution-model assessment are dated curation, not independently established conclusions in this draft.
\subsection{Short English statement}
For a fixed prime \(p\), the question asks how frequently the full Mordell--Weil--Selmer--Tate--Shafarevich exact sequence of an elliptic curve over \(k\) is isomorphic to a prescribed abstract exact sequence of \(\Z_p\)-modules. The field assumptions and the meaning of density are missing from the supplied statement, so the number-field formulation and counting limit above are conditional. This is a statement-only draft requiring human mathematical review.
\subsection{Sources}
\begin{itemize}
\item Supplied local UnsolvedMath registry record, permanent ID \(30002158\), title \emph{Distribution of Elliptic-Curve Selmer Exact Sequences}; source problem code \texttt{catalogue entry}. Repository registry locations are \texttt{data/unsolvedmath/problems.json} and the supplementary \texttt{catalog/problems.json}; the supplied record is the evidence used here.
\item Supplied report citation: \emph{Arithmetic Geometry (2013)}, Oberwolfach Reports, DOI \texttt{10.4171/owr/2012/38}, pp.\ 2335--2388, \url{https://doi.org/10.4171/owr/2012/38}. The record gives proposed year 2012 and supplies no individual proposer. Its verification note attributes the displayed sequence and question to Question 1.1 of the official report; that attribution is supplied information and has not been independently checked here.
\item Supplied literature source: Bhargava--Kane--Lenstra--Poonen--Rains, distribution model, \url{https://doi.org/10.4310/CJM.2015.v3.n3.a1}.
\item Supplied literature source: Oberwolfach report, \url{https://doi.org/10.4171/owr/2012/38}.
\item Literature assessment dated exactly \texttt{2026-08-21}: the supplied curation labels the problem open and states that the literature constructs a model distribution while its arithmetic identification is conjectural. This is supplied dated curation, not an independently verified literature assessment; its time-sensitive status and cited claims require independent review.
\end{itemize}
\end{document}
